\section{Практическая часть}

\subsection{Построение теоретических графиков зависимости $I(\theta)$}
 Постройте теоретические графики зависимости $I(\theta)$ для дифракции на одной, двух и 15 щелях, рассчитав (с точностью до $1''$) положение минимумов (нулей) для одной и двух щелей и главных максимумов для 15 щелей. Длину волны считайте равной 650 нм (красный свет), параметры щелей возьмите из таблицы:

\begin{table}[H]
	\centering
	\begin{tabular}{ccc}
		\toprule
		Число щелей & Ширина щели $b$ (мм) & Период решетки $d$ (мм) \\
		\midrule
		1  & 0,52 & -- \\
		2  & 0,52 & 1,5 \\
		15 & 1    & 2 \\
		\bottomrule
	\end{tabular}
\end{table}

Обратимся к теории, в частности к формуле с учётом, что $k = 2 \pi / \lambda$:
\[
I(\theta) = I_{max}\left( \frac{\sin \xi}{\xi} \right)^2, \quad \xi = \frac{\pi b}{\lambda} \sin \theta
\]

Найдём минимумы для $I(\theta)$. Очевидно, они будут совпадать с нулями числителя: $\sin \xi = 0$, что соответствует углам $\xi = \pi n, n \in \Z$, выразим угол $\theta$:
\[
\frac{\pi b}{\lambda} \sin \theta = \pi n 
\quad \Rightarrow \quad 
\sin \theta = \frac{\lambda}{b} n
, \quad n \in \Z
\]
\begin{align} \label{theta-theory}
	\theta = \arcsin{\frac{\lambda}{b}n}, \quad n \in \Z
\end{align}

Найдём углы для первых 6-ти нулей (с точностью до секунд) и запишем их в таблицу. Будем искать для положительных $n$, так как функция симметрична. Так же нулевой минимум совпадает с нулём знаменателя, то есть является главным максимумом. 

Так же учтём перевод из радиан в градусы:
\begin{align} \label{coef-angle}
	\theta^\circ = \frac{180}{\pi} \theta_{\text{рад}}
\end{align}

С учётом \cref{theta-theory} и \cref{coef-angle} итоговая формула для $\theta$ имеет вид:
\begin{align}
	\theta = \frac{180}{\pi} \arcsin{ \frac{\lambda}{b} n }
\end{align}


\renewcommand{\arraystretch}{1.8}
\begin{table}[H]
	\centering
	\caption{Углы минимумов для решётки с одной щелью}
	\begin{tabular}{cccc}
		\toprule
		Номер $n$ & $\sin \theta$ & $\theta$ & $\Delta \theta$ \\
		\midrule
		1 & $\dfrac{\lambda}{b}$  & $ 4' 18''$  & -- \\
		2 & $\dfrac{2\lambda}{b}$ & $ 8' 36''$ & $4' 18''$ \\
		3 & $\dfrac{3\lambda}{b}$ & $12' 54''$ & $4' 18''$ \\
		4 & $\dfrac{4\lambda}{b}$ & $17' 12''$ & $4' 18''$ \\
		5 & $\dfrac{5\lambda}{b}$ & $21' 30''$ & $4' 18''$ \\
		6 & $\dfrac{6\lambda}{b}$ & $25' 48''$ & $4' 18''$ \\
		\bottomrule
	\end{tabular}
\end{table}
\renewcommand{\arraystretch}{1}

Построим теоретический график $I(\theta)$, для одной щели(см. \cref{graph:1})

При дифракции плоской волны на решётке из $N$ щелей ширины $b$ с периодом $d$ области Фраунгофера соответствует распределение интенсивности:
\begin{align} \label{I-for-N}
	I(\theta) = I_{max}\left( \frac{\sin \xi}{\xi} \right)^2 \left( \frac{\sin{N\eta}}{\sin \eta} \right)^2, \quad \eta = \frac{kd}{2} \sin \theta
\end{align}

Заметим, что по формуле \cref{I-for-N} распределение интенсивности для решётки с одной щелью является огибающей графика(щелевой множитель) распределения интенсивности для решётки с $N$ щелями (решёточный множитель). 

Формула \cref{I-for-N} хорошо подходит для описания распределения с вспомогательными максимумами, но для решётки с одной щелью их нет, поэтому преобразуем вторую скобку(решёточный множитель) по формуле синуса двойного угла, чтобы найти его минимумы: 
\[
\frac{\sin N\eta}{\sin \eta} = 
\big[\sin 2x = 2 \sin x \cos x, N = 2\big] = 
\frac{2 \sin \eta \cos \eta }{ \sin \eta} =
2 \cos \eta
\]

Формула \cref{I-for-N} принимает минимальное значение, если косинус равен нулю, это происходит при, что $\eta = \pi / 2 + \pi m, m \in Z$. Найдём углы $\theta$ для минимумов решёточного множителя:
\begin{align}
	\frac{kd}{2} \sin \theta = \pi / 2 + \pi m \Rightarrow \theta = 
	\arcsin \big( \frac{\lambda}{2d}(2m + 1) \big), \quad m \in \Z
\end{align}

С учётом \cref{coef-angle} мы получим следующую формулу для углов минимумом решёточного множителя:
\[
\theta = \frac{180}{\pi}
\arcsin \big( \frac{\lambda}{2d}(2m + 1) \big), \quad m \in \Z
\]

Запишем полученные углы для первых 6-ти минимумов в таблицу:

\renewcommand{\arraystretch}{1.9}
\begin{table}[H]
	\centering
	\caption{Углы минимумов для решётки с двумя щелями}
	\begin{tabular}{cccc}
		\toprule
		Номер $m$ & $\sin \theta$ & $\theta$ & $\Delta \theta$ \\
		\midrule
		0 & $\dfrac{\lambda}{2d}$   & $ 0' 44''$        & -- \\
		1 & $\dfrac{3\lambda}{2d}$  & $ 2' 14''$ & $1' 30''$ \\
		2 & $\dfrac{3\lambda}{2d}$  & $ 3' 44''$ & $1' 30''$ \\
		3 & $\dfrac{5\lambda}{2d}$  & $ 5' 14''$ & $1' 30''$ \\
		4 & $\dfrac{7\lambda}{2d}$  & $ 6' 44''$ & $1' 30''$ \\
		5 & $\dfrac{9\lambda}{2d}$  & $ 8' 14''$ & $1' 30''$ \\
		6 & $\dfrac{11\lambda}{2d}$ & $ 9' 44''$ & $1' 30''$ \\
		\bottomrule
	\end{tabular}
\end{table}
\renewcommand{\arraystretch}{1}
Построим теоретический график $I(\theta)$, для двух щелей(см. \cref{graph:2})

Главные максимумы распределения \cref{I-for-N} определяются нулями знаменателя решёточного множителя $\sin \eta = 0$, найдём соответствующие им углы:
\[
\frac{kd}{2} \sin \theta = \pi s \Rightarrow \theta = \arcsin{\frac{\lambda}{d}s}, \quad s \in \Z
\]

Итоговая формула примет вид:
\[
\theta = \frac{180}{\pi} \arcsin{\frac{\lambda}{d}s}, \quad s \in \Z
\]

Найдём углы для первых двух максимумов и занесём в таблицу:
\renewcommand{\arraystretch}{1.8}
\begin{table}[H]
	\centering
	\caption{Углы главных максимумов для решётки с 15 щелями}
	\begin{tabular}{cccc}
		\toprule
		Номер $s$ & $\sin \theta$ & $\theta$ & $\Delta \theta$ \\
		\midrule
		1 & $\dfrac{\lambda}{d}$  & $1' 7''$ & -- \\
		2 & $\dfrac{2\lambda}{d}$ & $2' 14''$ & $1' 7''$ \\
		\bottomrule
	\end{tabular}
\end{table}
\renewcommand{\arraystretch}{1}

Построим теоретический график $I(\theta)$, для 15-ти щелей(см. \cref{graph:3})

\subsection{Измерение углов}

Используя ГС-30 найдём углы на которых наблюдаются минимумы интенсивности(чёрные полосы). 

Методика измерения будет следующей: наводим перекрестие зрительной трубки на середину самой правой измеряемой полосы, записываем показания гониометра, далее перемещаем перекрестие на соседнюю полосу слева. Повторяем до самого левого измеряемого минимума. Полученные значения занесём в таблицу.

%\renewcommand{\arraystretch}{1.8}
\begin{table}[H]
	\centering
	\caption{Углы минимумов для решётки с одной щелью}
	\begin{tabular}{ccc}
		\toprule
		Номер & $\theta$ & $\Delta \theta$ \\
		\midrule
		-3 & $ 185^\circ  7' 30''$  & -- \\
		-2 & $ 185^\circ  3' 15''$ & $4' 15''$ \\
		-1 & $ 184^\circ 59' 15''$ & $4' 00''$ \\
		 1 & $ 184^\circ 51' 00''$ & $4' 15''$ \\
		 2 & $ 184^\circ 46' 45''$ & $4' 15''$ \\
		 3 & $ 184^\circ 42' 30''$ & $4' 15''$ \\
		\bottomrule
	\end{tabular}
\end{table}
%\renewcommand{\arraystretch}{1}

Измерьте углы, под которыми наблюдаются минимумы интенсивности при дифракции на двух щелях.
\begin{table}[H]
	\centering
	\caption{Углы минимумов для решётки с двумя щелями}
	\begin{tabular}{ccc}
		\toprule
		Номер & $\theta$ & $\Delta \theta$ \\
		\midrule
		-6 & $185^\circ 1' 14''$  & -- \\
		-5 & $185^\circ 0' 15''$  & $1' 30''$ \\
		-4 & $184^\circ 59' 30''$ & $0' 45''$ \\
		-3 & $185^\circ 58' 45''$ & $0' 45''$ \\
		-2 & $185^\circ 57' 15''$ & $1' 30''$ \\
		-1 & $184^\circ 56' 45''$ & $1' 30''$ \\
		 1 & $184^\circ 54' 15''$ & $1' 30''$ \\
		 2 & $184^\circ 52' 45''$ & $1' 30''$ \\
		 3 & $184^\circ 51' 0''$  & $1' 45''$ \\
		 4 & $184^\circ 50' 15''$ & $0' 45''$ \\
		 5 & $184^\circ 49' 30''$ & $0' 45''$ \\
		 6 & $184^\circ 48' 0''$  & $1' 30''$ \\
		\bottomrule
	\end{tabular}
\end{table}

Для максимумов мы просто измеряем расстояние между яркими полосами (соответственно максимумами).
Измерим углы, под которыми наблюдаются главные максимумы интенсивности при дифракции на решётке из 15 щелей.

\begin{table}[H]
	\centering
	\caption{Углы максимумов для решётки с 15 щелями}
	\begin{tabular}{ccc}
		\toprule
		Номер & $\theta$ & $\Delta \theta$ \\
		\midrule
		-1 & $ 184^\circ 56' 0''$  & -- \\
		0  & $ 184^\circ 54' 45''$ & $1' 15''$ \\
		1  & $ 184^\circ 53' 45''$ & $1' 0''$ \\
		\bottomrule
	\end{tabular}
\end{table}

Сравним измеренные в заданиях 3--5 углы с результатами теоретических расчётов в задании 1.

\subsection{Определение размера источника света при дифракции на двух щелях}
Наблюдая дифракцию на двух щелях, определите размер источника света (ширину входной щели коллиматора), при котором наступает размытие дифракционной картины, и проверьте соотношение (7).
